\section{Scope, conventions, and the three meanings of BCH}\label{sec:scope}
The equation
\begin{equation}\label{eq:main}
 e^Xe^Y=e^Z
\end{equation}
is not, by itself, a definition of a unique logarithm. The exponential can fail to be injective, and a product of exponentials need not be an exponential \emph{in the specified Lie algebra}. What is universal is the formal series
\begin{equation}\label{eq:formal-bch}
 Z(X,Y)=\BCH(X,Y)=\log(e^Xe^Y)=\sum_{n\geq1}Z_n(X,Y),
\end{equation}
where $Z_n$ has total degree $n$. A principal matrix logarithm, a continued branch of a logarithm, and a convergent BCH series are related but distinct objects.

The coverage target is the version of the source article identified by revision number 1368909113, consulted on 16 September 2026 \cite{Wikipedia}. Every displayed mathematical identity and every mathematical assertion stated in its main text or mathematical notes is addressed below. Historical priority statements, application surveys, and theorems merely named in links are not additional proof obligations; in particular, the article does not state the Stone--von Neumann or Golden--Thompson theorems. \Cref{sec:coverage} gives a claim-by-claim concordance. Several assertions about geometry require correction, not proof as written: $\ad_X$ is generally singular and is not the coframe matrix; a Killing form need not be a metric; and an arbitrary pullback can be degenerate. We give the correct statements and explicit counterexamples.

\subsection{Algebraic and analytic settings}
Throughout the formal arguments, $\kk$ is a field of characteristic zero. Let
\[
 \TT=\kk\langle X,Y\rangle=\bigoplus_{n\geq0}\TT_n,
 \qquad \widehat\TT=\prod_{n\geq0}\TT_n,
\]
where $\TT_n$ has as a basis the words of length $n$ in $X,Y$. The empty word is the unit $1$. Multiplication is concatenation, extended degree by degree. Thus the coefficient of a word in a product is always a finite sum. Write
\[
 F^r\widehat\TT=\prod_{n\geq r}\TT_n.
\]
A sum converges \emph{formally} when every fixed degree eventually stabilizes. No norm or numerical limiting process is involved. The bracket is
\begin{equation}\label{eq:bracket}
 [A,B]=AB-BA,\qquad \ad_A B=[A,B].
\end{equation}
Associativity gives antisymmetry and the Jacobi identity by expansion of the twelve resulting products. In particular,
\begin{equation}\label{eq:derivation}
 [A,[B,C]]=[[A,B],C]+[B,[A,C]],
 \qquad [\ad_A,\ad_B]=\ad_{[A,B]}.
\end{equation}
The Lie subalgebra of $\TT$ generated by $X,Y$ is denoted $\LL(X,Y)$; its degree completion is $\widehat\LL(X,Y)$.

For analytic statements, $\Algebra$ is a real or complex unital Banach algebra with a submultiplicative norm. Finite matrices are the main example. When discussing arbitrary finite-dimensional Lie groups we use their ordinary smooth exponential map and derive the same formulas intrinsically. Unbounded Hilbert-space operators are handled separately in \cref{sec:quantum}.

\begin{lemma}[Formal exponential and logarithm]\label{lem:exp-log}
The series
\begin{equation}\label{eq:exp-log}
 \exp A=\sum_{n\geq0}\frac{A^n}{n!},\qquad
 \log(1+U)=\sum_{n\geq1}\frac{(-1)^{n-1}}{n}U^n
\end{equation}
are mutually inverse bijections between $F^1\widehat\TT$ and $1+F^1\widehat\TT$.
If $A,B$ commute, then $e^{A+B}=e^Ae^B$; if $g,h\in1+F^1\widehat\TT$ commute, then $\log(gh)=\log g+\log h$.
\end{lemma}
\begin{proof}
Positive order makes every coefficient in each substitution a finite sum. The scalar identities $\log(\exp t)=t$ and $\exp(\log(1+t))=1+t$ follow, for example, by differentiating the corresponding formal scalar series and checking the constant terms. Substitution of the single element $A$ or $U$ is legitimate: its powers commute with each other even though the ambient algebra need not be commutative. For commuting $A,B$, the binomial theorem gives
\[
 \sum_{n\geq0}\frac{(A+B)^n}{n!}
 =\sum_{r,s\geq0}\frac{A^rB^s}{r!s!}=e^Ae^B.
\]
When $g,h$ commute, their logarithms commute as well, so the preceding identity and the inverse property give the last assertion.
\end{proof}

In a Banach algebra the exponential series converges for every $A$, since
$\norm{A^n/n!}\leq\norm A^n/n!$. The logarithm series converges whenever $\norm U<1$. Formal identities involving absolutely convergent substitutions remain valid analytically.

\subsection{How a universal series is evaluated}
A homogeneous Lie polynomial can be evaluated at any pair of elements of any characteristic-zero Lie algebra. The identification of $\LL(X,Y)$ with the abstract free Lie algebra, and the fact that the free associative algebra is its universal enveloping algebra, are proved in \cref{sec:pbw}. For a general Lie algebra without a topology, the unambiguous statement is
\[
 \BCH(tX,tY)=\sum_{n\geq1}t^n Z_n(X,Y)\in t\gfrak[[t]],
\]
not that the sum at $t=1$ is an element of $\gfrak$. In a complete positively filtered Lie algebra it can instead be evaluated by filtration convergence. In a nilpotent Lie algebra it is a finite polynomial.

Our convention for an iterated commutator is \emph{right nesting}:
\begin{equation}\label{eq:right-word}
 R(a_1\cdots a_n)=[a_1,[a_2,\ldots,[a_{n-1},a_n]\ldots]],
 \qquad R(a_1)=a_1.
\end{equation}
The letters may be $X$ or $Y$. Thus $R(XXY)=\ad_X^2Y$ and $R(YYX)=\ad_Y^2X$. The symbols $X^rY^s$ inside $R$ describe a word, not a commutator of associative powers.
